
\subsubsection{2.1 Uncertainty Quantification in LLMs}

Early work on LLM uncertainty relies on token-level conditional probabilities, but RLHF-trained models exhibit poor calibration between confidence and accuracy. Verbalized confidence often outperforms raw probabilities for post-RLHF models.

Semantic entropy addresses linguistic invariance by clustering semantically equivalent generations before computing entropy, achieving AUROC approximately 0.75 on TruthfulQA but requiring 5-10 samples per query. SelfCheckGPT provides black-box consistency checking via sampling and cross-referencing, again requiring multiple generations.

Probe-based methods avoid multi-sampling by training classifiers on internal representations. MIND uses internal states for real-time unsupervised detection. Semantic entropy probes demonstrate that lightweight linear classifiers on hidden states approach the accuracy of full semantic entropy computation.

\subsubsection{2.2 Cross-Layer Analysis}

The logit lens projects intermediate representations through the final unembedding matrix, revealing per-layer probability distributions. END decoding demonstrates that cross-layer entropy correlates with factuality. The present work formalizes this observation into trajectory metrics (NTI, CMI) and validates them with controlled experiments.

\subsubsection{2.3 Positioning}

\begin{table}[htbp]
\centering
\resizebox{\columnwidth}{!}{%
\begin{tabular}{llll}
\toprule
Method & Single-Pass & No External DB & Interpretable \\
\midrule
Semantic Entropy & No & Yes & Moderate \\
SelfCheckGPT & No & Yes & High \\
MIND & Yes & Yes & Low \\
CLTI (this work) & Yes & Yes & High \\
\bottomrule
\end{tabular}
}
\end{table}

